\section{Observation separation and lower bounds}
\label{sec:negative-proofs}

Throughout this section a word has symbols $X_0,\ldots,X_{M+1}$,
pair states $Y_i=(X_i,X_{i+1})$, and departures $0\le i<M$.
Write $J_{ab}=\sum_{i<M}\ind\{Y_i=ab\}$,
$J_{\rm cross}=J_{01}+J_{10}$, and $N_{abc}=H_{abc}(X)$.
Thus $Y_M$ is a terminal pair, not a departure. All risks below include
independent predictor randomness. The observation $\mathcal T_r$ consists
of the eight triple counts and the actual first and last $r$ symbols.

\subsection{An exact occupation mixture}

For $0<a,t<1$, let $Q_a$ be the stationary law with pair transition matrix,
in the order $00,01,10,11$,
\[
 P_a=\begin{pmatrix}
 1-a&a&0&0\\
 0&0&1/2&1/2\\
 1/2&1/2&0&0\\
 0&0&t&1-t
 \end{pmatrix},\qquad
 \Pi_a=\frac{(t/a,2t,2t,1)}{1+t/a+4t}.
\]
In particular the stationary initial factor is
$\Pi_a(11)=a/[t+(1+4t)a]$.

\begin{lemma}[Occupation-mixture domination]
\label{lem:mixture}
Put $F=\{Y_0=11,Y_M=00\}$ and $U=N_{001}$. For every integer $j$,
every integer $u_0\ge1$, every interval $I\subset(0,1)$, and every
additional event $G$ (possibly depending on $a$),
\[
 \int_I Q_a(F,J_{00}=j,U\ge u_0,G)\,\frac{da}{a}
 \le \frac{t}{u_0}.
\]
\end{lemma}
\begin{proof}
Delete every transition departing from $00$, both loops and exits.
For $J_{00}=j$, exactly $L=M-j$ transitions remain. Their last transition
is $10\to00$. Before that transition the retained path is on
$B=01,C=10,D=11$. A retained $C\to00$ followed by an erased exit
$00\to B$ is represented by a marked $C\to B$ of weight $1/2$;
a direct $C\to B$ is unmarked and also has weight $1/2$. The erased
exit's factor $a$ is restored separately. All other retained transitions
keep their probabilities. The marked transition matrix is therefore
\[
 K(z)=\begin{pmatrix}
 0&1/2&1/2\\ (1+z)/2&0&0\\ 0&t&1-t
 \end{pmatrix}.
\]
For $L\ge1$, define
\[
 h_{L,u}=\tfrac12[z^u](K(z)^{L-1})_{D,C},
\]
and set $h_{L,u}=0$ for $L\le0$ or $u<0$. The final transition
$C\to00$ has weight $1/2$ and is not marked. A skeleton with $u$ marks
has $u$ completed visits to $00$ and one terminal visit. Distributing
its $j-u$ loops among these $u+1$ holding slots gives exactly
$\binom ju$ weak compositions. Each composition and marked skeleton
recover one word, and deletion recovers both uniquely. Consequently,
for $0\le u\le j$,
\begin{equation}
 Q_a(F,J_{00}=j,U=u)
 =\Pi_a(11)\binom ju a^u(1-a)^{j-u}h_{M-j,u}.
 \label{eq:marked-word-identity}
\end{equation}
Infeasible cases have zero mass. The terminal holding slot is present
even though $Y_M$ is not a departure; there is no initial $00$ slot.

For $u\ge1$, the beta integral gives
$\binom ju\mathrm B(u,j-u+1)=1/u$. Thus, if
$A_{j,u}$ has the $\operatorname{Beta}(u,j-u+1)$ law,
\begin{align}
 \int_0^1 Q_a(F,J_{00}=j,U=u)\,\frac{da}{a}
 &=\frac{h_{M-j,u}}{u}\E\Pi_{A_{j,u}}(11)
 \le\frac{h_{M-j,u}}{u}. \label{eq:stationary-beta-factor}
\end{align}
The expectation is part of the equality: stationarity has not been
cancelled. For $u=0$ the corresponding factor is instead
$h_{M-j,0}\int_0^1(1-a)^j/[t+(1+4t)a]\,da$, and no $1/u$ formula
is used. This is why $u_0>0$ is required.

To sum over marks at the exact outside horizon, put $R=K(1)$.
This matrix is stochastic. Started at $D$, let $b_n,c_n,d_n$ be its
probabilities of $B,C,D$. Then
\[
 c_0=0,\quad c_1=t,\quad b_n=c_{n-1}\ (n\ge1),\qquad
 c_n=\tfrac12 c_{n-2}+td_{n-1}\ (n\ge2).
\]
Induction gives $c_n\le2t$. Hence, for each $L\ge1$,
\[
 \sum_{u\ge0}h_{L,u}=\tfrac12 c_{L-1}\le t.
\]
First drop $G$ and enlarge $I$ to $(0,1)$ in the nonnegative integral.
Then sum \eqref{eq:stationary-beta-factor} over $u\ge u_0$ to obtain
\[
 \int_I Q_a(F,J_{00}=j,U\ge u_0,G)\,\frac{da}{a}
 \le\sum_{u\ge u_0}\frac{h_{M-j,u}}u
 \le\frac1{u_0}\sum_{u\ge0}h_{M-j,u}\le\frac t{u_0}.
\]
No conditioning on a prescribed number of outside marks or replacement
of the horizon $M-j$ is involved.
\end{proof}

\subsection{Relocating loops while preserving sampled windows}

\begin{lemma}[Interior loop relocation]
\label{lem:relocation}
Suppose $Y_0=11,Y_M=00$, $2\le r$, and $2r\le M+2$.
Put $l=r-2$, $h=M-r+2$. If the interior vertices
$Y_l,\ldots,Y_h$ contain a visit to $11$ strictly before a visit to $00$,
there is a deterministic transformation $C_r$ preserving the complete
triple histogram and the actual first and last $r$ symbols, such that,
with $j=J_{00}$,
\begin{equation}
 d_H\bigl(C_r(X),1^{M+2-j}0^j\bigr)
 \le3J_{\rm cross}+4r-5.
 \label{eq:relocation-distance}
\end{equation}
If at least one of the two movable diagonal-loop types is absent, the
same conclusion holds without the ordered-visit condition.
\end{lemma}
\begin{proof}
Freeze departures $i<l$ and $i\ge h$. These are $2r-4$ edges and
spell exactly the sampled boundary symbols. From the interior path
remove only the $111$ and $000$ loops, in numbers $A_1$ and $A_0$,
retaining every other edge in order. Let $S_0,\ldots,S_L$ be its vertices.
Deleting self-loops neither removes a distinct-state visit nor reverses
its order. Choose the first retained $11$ and the first subsequent
retained $00$; insert all $A_1$ loops at the former and all $A_0$ at
the latter. The choice is deterministic. Deleting or inserting a loop
preserves its attachment pair. Reinserting exactly the removed number
of edges preserves the absolute suffix attachment index $h$, so the
frozen edge labels and all sampled boundary symbols are unchanged.
Every triple label is retained or reinserted once, proving histogram
preservation. Unfrozen boundary-crossing triples need not keep their
individual positions.

Here is the required positional count argument. A binary word of length
$N$ with a distinguished block of $A_1$ ones before a distinguished
block of $A_0$ zeros has $o$ other ones and $z$ other zeros. If
$j=A_0+q$, $q\ge0$, then
\begin{equation}
 d_H(\text{word},1^{N-j}0^j)
 \le\max\{z+q,\,2o+z-q\}. \label{eq:two-block-distance}
\end{equation}
Indeed, write $z_-$ for zeros before the target cut. The exact mismatch
count is $2z_-+j-(A_0+z)$. If a distinguished one lies after the cut,
all distinguished zeros lie after the cut and $z_-\le z$, giving
$z+q$. Otherwise every one after the cut is among the $o$ other ones;
writing $o_+$ for their number, the other exact identity is
$2o_++(A_0+z)-j$, at most $2o+z-q$. This proof does not presume any
order among the remaining symbols.

Put $u=N_{001}$, $v=N_{011}$, $p=N_{010}$, $w=N_{101}$, and
$d=J_{\rm cross}$. Endpoint flow gives
\[
 N_{100}=u+1,\qquad N_{110}=v+1,\qquad
 d=p+v+u+1+w.
\]
Let $f_1,f_0$ be the frozen $111,000$ loop counts, so
$f_1+f_0\le2r-4$. The initial pair contributes two other ones.
Counting the appended symbol of every remaining edge gives exactly
\begin{align*}
 o&=2+u+v+w+f_1=d+1-p+f_1,\\
 z&=p+u+v+2+f_0=d+1-w+f_0,\\
 q&=j-A_0=u+f_0\ge0.
\end{align*}
Since $u\le d-1$ and $d\ge1$,
\begin{align*}
 z+q&=d+1-w+u+2f_0\le2d+2f_0,\\
 2o+z-q&=3d+3-2p-w-u+2f_1\le3d+3+2f_1.
\end{align*}
Each is at most $3d+4r-5$, proving \eqref{eq:relocation-distance}.
If only one movable loop type occurs, put its block at any retained
visit to its state, which exists because the loops were deleted there.
The same two-block argument treats the absent block as empty; if both
are absent use the identity transformation.

Finally, $j$ is the sampled word's actual departure count. The comparison
word has instead
\begin{equation}
 J_{00}(1^{M+2-j}0^j)=j-2\quad(2\le j\le M).
 \label{eq:null-suffix-shift}
\end{equation}
For this range its final pair and the transformed word's final pair are
both $00$, so all their Hamming differences are in the editable first
$M$ coordinates.
\end{proof}

\subsection{A finite matched-observation lower bound}

\begin{proposition}[Finite critical lower bound]
\label{prop:finite-lower}
Let
\[
 c=2^{-16},\quad A=2^{12},\quad t=\frac{ck}{M},\quad
 u_0=\lfloor ck/4\rfloor,\quad T=\lfloor M/8\rfloor,\quad
 L_-=\log\frac1{4At}.
\]
Assume integer $M,k,r$ satisfy
\begin{equation}
 \begin{gathered}
 M\ge64,\quad 2\le r\le k/128,\quad k\le M/8,\quad ck\ge1024,\\
 At<1/4,\qquad \frac M{k^2}\le\frac{Ac^2}{2^{20}}.
 \end{gathered}\label{eq:finite-lower-domain}
\end{equation}
For any fixed predictor based on $\mathcal T_r$, let $S$ be its supremum
clean risk over the stationary laws $Q_a$, $At\le a\le1/4$.
Every such source satisfies the all-state certificate with supplied
$\beta=1/M$. There is a stationary source with the same supplied
certificate and one source-independent strict-$k$ corruption kernel
for which
\begin{equation}
 R_{\rm null}\ge\frac{u_0}{8Mt}
       \left[\left(\frac t{16}-S\right)L_-\right]_+.
 \label{eq:finite-lower}
\end{equation}
The kernel protects the final pair on every input and every output.
\end{proposition}
\begin{proof}
We give the source, probability, and kernel calculations separately.

\paragraph{Source legality and all sixteen certificate entries.}
Throughout $t\le c/8$, $a\ge At\ge t$, and $a\le1/4$. Set
$D=1+t/a+4t$. Then
$1\le D\le1+1/A+c/2<4/3<2$.
The flow equalities
$\Pi_a(00)a=\Pi_a(10)/2$ and
$\Pi_a(11)t=\Pi_a(01)/2$, together with
$\Pi_a(01)=\Pi_a(10)$, verify stationarity directly. Multiplication gives
\begin{equation}
 P_a^2=\begin{pmatrix}
 (1-a)^2&(1-a)a&a/2&a/2\\
 1/4&1/4&t/2&(1-t)/2\\
 (1-a)/2&a/2&1/4&1/4\\
 t/2&t/2&(1-t)t&(1-t)^2
 \end{pmatrix}.
 \label{eq:all-sixteen}
\end{equation}
Every displayed entry is at least $t/2$: use $a\ge t$ and
$a,t\le1/4$ for the products and all other entries are immediate.
Since $t/2=ck/(2M)\ge1/M$, every starting and destination state obeys
$P_a^2(u,v)\ge M^{-1}\Pi_a(v)$. Also $a\ge At\ge2/M$, so
$a-\beta\ge a/2$.

\paragraph{Common-atom covariance and coupling bounds.}
The four paths to $11$ of length three, from the states in order, have
probabilities
\[
 (1-a)a/2,\qquad (1-t)^2/2,\qquad a/4,\qquad(1-t)^3.
\]
These arise respectively from paths $00,00,01,11$;
$01,11,11,11$; $10,00,01,11$; and $11,11,11,11$.
Each is at least $\lambda=a/4$, so
$P_a^3(x,\cdot)\ge\lambda\delta_{11}$ for every $x$.
Split each three-step transition into this common atom and a residual
probability law. Shared independent Bernoulli trials of probability
$\lambda$ couple any two chains at the first common atom. Fill in each
three-step block by its conditional bridge; this preserves the full
one-step path law. After meeting, use identical transitions. Thus the
contraction coefficient at lag $h$ is bounded by
\[
 \alpha_h=(1-\lambda)^{\lfloor h/3\rfloor}.
\]

For a stationary chain $W$, a departure-state-set count with stationary
mass $p$, or a count of a specified directed edge with stationary
probability $p$, satisfies
\begin{equation}
 \operatorname{Var}(V)\le\frac{7Mp}{\lambda}.
 \label{eq:common-atom-variance}
\end{equation}
For an edge indicator $I_i$ and $h\ge1$, condition on the state after
edge $i$. The conditional probability of edge $i+h$ is a $[0,1]$-valued
function after $h-1$ transitions, whose deviation from its stationary
value is at most $\alpha_{h-1}$. Therefore
$|\operatorname{Cov}(I_i,I_{i+h})|\le p\alpha_{h-1}$.
For a state-set indicator use $\alpha_h\le\alpha_{h-1}$ instead.
Since $\sum_{h\ge1}\alpha_{h-1}=3/\lambda$, expansion of the variance
gives $Mp(1+6/\lambda)\le7Mp/\lambda$, as claimed.

Couple a chain $Z$ started at $00$ to stationary $W$ by the same block
construction. Its meeting boundary $\sigma$ has
\[
 \E\sigma\le3/\lambda=12/a.
\]
Every departure-state or edge count in the two paths differs by at most
$\sigma$ pathwise: after boundary $\sigma$ the paths agree, and there
are at most $\sigma$ earlier departures. No independence of $\sigma$
and the stationary counts is needed.

\paragraph{Stationary reversal and count corrections.}
Let $\mathcal R(ab)=ba$. The stationary triple identity
\[
 \Pi_a(ab)P_a(ab,bc)=\Pi_a(cb)P_a(cb,ba)
\]
holds for every triple: the only nonpalindromic cases are $001,100$
and $011,110$, given by the two flow identities above. It implies
$\mathcal R P_a^*\mathcal R=P_a$, where $P_a^*$ is the stationary
reverse pair kernel. Reversing the symbols of a stationary word thus
preserves its law. In particular, with
$F=\{Y_0=11,Y_M=00\}$ and $E=\{Z_M=11\}$,
\[
 Q_a(F)=\Pi_a(00)P_{00}(E),
\]
and conditional on $F$, the reversed pair path
$Z_i=\mathcal RY_{M-i}$ has law $P_{00}(\cdot\mid E)$.
For $M\ge64$, $\lfloor M/3\rfloor\ge M/4$, and therefore
\[
 \alpha_M\le e^{-aM/16}\le\frac{16}{aM}
 \le\frac{16}{Ack}\le\frac14.
\]
Since $\Pi_a(11)=1/D\ge3/4$, we have $P_{00}(E)\ge1/2$.

Put $V(W)=N_{001}(W)$ and $x=ck$. Its stationary mean is
$Mt/D\in[x/2,x]$. Consider the four conditions
\begin{equation}
 \sigma\le x/16,\quad |V(W)-Mt/D|\le x/8,\quad
 J_{\rm cross}(W)\le k/32,\quad J_{00}(W)\le M/16.
 \label{eq:preliminary-good}
\end{equation}
Their failure probabilities, respectively, are at most
\begin{equation}
 \frac{192M}{Ac^2k^2},\quad
 \frac{1792M}{Ac^2k^2},\quad
 \frac{458752M}{Ak^2},\quad \frac{16}{A}.
 \label{eq:four-failures}
\end{equation}
The first follows from $\E\sigma\le12/a$ and $a\ge Ack/M$.
For the second, \eqref{eq:common-atom-variance} with $p=t/D$ gives
$\operatorname{Var}V\le28Mt/a$; divide by $(ck/8)^2$.
For the third, the mean is $4Mt/D\le4ck\le k/64$, and
$\operatorname{Var}J_{\rm cross}\le112Mt/a\le112M/A$;
divide by $(k/64)^2$. For the last use
$\Pi_a(00)=t/(aD)\le1/A$ and Markov's inequality.
Conditioning the joint coupling on $E$ costs at most two. Under
\eqref{eq:finite-lower-domain}, twice the sum in
\eqref{eq:four-failures} is bounded by
\[
 2\left(\frac{1984}{2^{20}}
       +\frac{458752c^2}{2^{20}}+\frac{16}{A}\right)<\frac1{16}.
\]

On $F$, the exact reversal identities are
\begin{equation}
 \begin{split}
 J_{00}(Z)&=J_{00}(Y)+1,\\
 J_{\rm cross}(Z)&=J_{\rm cross}(Y),\\
 N_{001}(Z)&=N_{100}(Y)=N_{001}(Y)+1.
 \end{split}\label{eq:reverse-count-shifts}
\end{equation}
The first compares arrival indices $1,\ldots,M$ with departure indices
$0,\ldots,M-1$; the second loses no mixed endpoint; the third combines
triple reversal with entry-minus-exit flow at $00$.
Thus \eqref{eq:preliminary-good}, the pathwise coupling bound, and
$x\ge1024$ imply for the original path
\begin{align*}
 U&\ge x/2-x/8-x/16-1\ge x/4\ge u_0,\\
 J_{\rm cross}&\le k/32+ck/16\le k/16,\\
 2\le u_0&\le J_{00}\le M/16+ck/16\le M/8.
\end{align*}
The integer upper count bound is $J_{00}\le T$.

\paragraph{Ordered visits in the interior.}
Let $l=r-2$ and $h=M-r+2$. Since $r\le M/1024$,
$h-l=M-2r+4\ge M/2$. Write $B$ for a visit to $11$ before a later
visit to $00$ within $Y_l,\ldots,Y_h$. Because $M-h=l$ and $M-l=h$,
reversal sends this event precisely to a visit to $00$ before a later
visit to $11$ within $Z_l,\ldots,Z_h$.

The ordinary expected hitting times of $00$ and $11$, zero at their
targets and listed in state order, are
\begin{align*}
 h_{00}&=(0,\,4+1/t,\,3+1/(2t),\,3+3/(2t)),\\
 h_{11}&=(3+3/(2a),\,3+1/(2a),\,4+1/a,\,0).
\end{align*}
For completeness, the first vector follows by solving
$h_{01}=1+(h_{10}+h_{11})/2$,
$h_{10}=1+h_{01}/2$, and
$h_{11}=1+th_{10}+(1-t)h_{11}$, with $h_{00}=0$.
For the second solve
$h_{00}=1+(1-a)h_{00}+ah_{01}$,
$h_{01}=1+h_{10}/2$, and
$h_{10}=1+(h_{00}+h_{01})/2$, with $h_{11}=0$.
On $a,t\le1/4$ the two maxima are
$3+3/(2t)$ and $3+3/(2a)$. From any state at time $l$, hitting $00$
and then $11$ therefore has expected additional time at most
\[
 6+\frac{3}{2t}+\frac{3}{2a}\le\frac6t.
\]
This uses the strong Markov property at the first hit: partition by its
finite time and state, apply the usual Markov property on each event,
and sum. Distinct targets make the ordering strict. Markov's inequality
over $h-l\ge M/2$ bounds the unconditioned failure probability by
$12/(tM)$. Dividing by $P_{00}(E)\ge1/2$ gives
\[
 Q_a(B^c\mid F)\le\frac{24}{ck}\le\frac3{128}.
\]
This estimate is uniform in the time-$l$ state; it assumes no
independence from the final state or the counts.

Let $G$ be the intersection of $B$ with
$U\ge u_0$, $J_{\rm cross}\le k/16$, and $2\le J_{00}\le T$.
The preceding count and hitting bounds give
$Q_a(G\mid F)\ge1-1/16-3/128>1/2$. Retaining the rare endpoint mass,
\begin{equation}
 Q_a(F\cap G)\ge\Pi_a(00)/4,\qquad
 (a-\beta)Q_a(F\cap G)\ge t/16.
 \label{eq:good-event-mass}
\end{equation}
Indeed $P_{00}(E)\ge1/2$, $a-\beta\ge a/2$, and
$a\Pi_a(00)=t/D\ge t/2$.

\paragraph{The null and its full word masses.}
Take a stationary first-order binary source with
$T_*(1,0)=q_*=1/(2M)$ and $T_*(0,1)=\beta=1/M$.
Its symbol masses are $\pi_*(0)=1/3$, $\pi_*(1)=2/3$.
Its stationary pair law is $\Pi_*(s,z)=\pi_*(s)T_*(s,z)$, and
for all starting and destination pair states,
\[
 \frac{P_*^2((u,v),(s,z))}{\Pi_*(s,z)}
 =\frac{T_*(v,s)}{\pi_*(s)}.
\]
For $(v,s)=(0,0),(0,1),(1,0),(1,1)$ these four ratios are
$3(1-\beta),\,3\beta/2,\,3q_*,\,3(1-q_*)/2$, respectively.
They are all at least $\beta+q_*\ge\beta$, checking every pair-state
certificate entry. The target at terminal $00$ is Bernoulli$(\beta)$,
not Bernoulli$(q_*)$.
For $2\le j\le T$ put $x^{(j)}=1^{M+2-j}0^j$.
These disjoint words have full stationary masses
\begin{equation}
 v_j=\frac23(1-q_*)^{M+1-j}q_*(1-\beta)^{j-1}
 \ge\frac1{8M}. \label{eq:null-word-masses}
\end{equation}
Bernoulli's inequality gives the prefix factor at least $1/2$ and
the suffix factor at least $7/8$, so in fact the product is at least
$7/(48M)>1/(8M)$. These masses include the initial symbol and every
one of the $M+1$ transitions.

\paragraph{A predeclared strict kernel with identity rejection.}
For $2\le j\le T$, define a finite measure on words by
\[
 \nu_j(dy)=\int_{At}^{1/4}
       Q_a(dy;F,U\ge u_0,J_{00}=j)\,\frac{da}{a},
 \qquad \mu_j=\nu_j(\text{all words}).
\]
Lemma~\ref{lem:mixture} gives $\mu_j\le t/u_0$.
For $\mu_j>0$, normalization $\nu_j/\mu_j$ is a probability law.
These laws depend only on declared $M,k,r$ and the auxiliary family,
not the true source, its parameters, its initial law, or a predictor seed.

On input $x^{(j)}$ with $\mu_j>0$, draw $Y$ once from
$\nu_j/\mu_j$. If $B$ holds, form the deterministic $C_r(Y)$ of
Lemma~\ref{lem:relocation}. Accept it only if its actual Hamming
distance from $x^{(j)}$ is at most $k$ and its final pair equals that
of $x^{(j)}$. Otherwise output the unchanged input. On zero-mass labels
and all other input words, also output the unchanged input. This defines
one source-independent kernel on all words, with a strict pointwise
budget and protected final pair, not merely an average budget.
It uses one draw, not repeated sampling until acceptance.

On $G$, Lemma~\ref{lem:relocation} bounds the distance by
$3k/16+4r\le3k/16+k/32<k$, so every check succeeds, and
\[
 \mathcal T_r(C_r(Y))=\mathcal T_r(Y).
\]
The same statistic identity holds on every transformed acceptance;
it is not asserted on identity rejections. These are the sampled
alternative's boundary windows, not necessarily the null input's
windows. Input-window edits are allowed and included in the Hamming
check. The suffix label obeys exactly
$J_{00}(x^{(j)})=j-2$ and $J_{00}(Y)=j$.

\paragraph{Nonnegative-loss ordering.}
For a sampled word $y$, let $\ell_a(y)$ be the predictor's expected
absolute error against $a$ on $\mathcal T_r(y)$, and let $\ell_*(y)$
be its corresponding error against $\beta$. On $F$ the clean actual
terminal target is $a$; the triangle inequality, also after averaging
independent decoder randomness, gives
$\ell_*(y)\ge a-\beta-\ell_a(y)$.
For $\mu_j>0$, \eqref{eq:null-word-masses} gives
\[
 \frac{v_j}{\mu_j}\ge\rho:=\frac{u_0}{8Mt}.
\]
First discard nonnegative losses from other null inputs and rejected
proposals. On the remaining nonnegative accepted losses, replace
$v_j/\mu_j$ by $\rho$. Next restrict to $G$, where acceptance and
statistic matching hold. Only then use the signed triangle bound:
\begin{align*}
 R_{\rm null}
 &\ge\rho\int_{At}^{1/4}
          \E_{Q_a}[\ind_{F\cap G}\ell_*(Y)]\,\frac{da}{a}\\
 &\ge\rho\int_{At}^{1/4}
    \left((a-\beta)Q_a(F\cap G)
          -\E_{Q_a}[\ind_{F\cap G}\ell_a(Y)]\right)\frac{da}{a}\\
 &\ge\rho\left(\frac t{16}-S\right)L_-.
\end{align*}
The restricted alternative loss is bounded by its unconditional clean
risk, hence by $S$. Risk is nonnegative, giving \eqref{eq:finite-lower}.
In particular, no lower mass multiplier has been placed in front of a
possibly negative quantity, and no endpoint-conditioned loss has been
substituted for the unconditional target risk.
\end{proof}

\subsection{Consequences for widths and supplemental messages}

\begin{proposition}[Width and fixed-fraction consequences]
\label{prop:lower-consequences}
For a prescribed sequence $2\le r_M\le M+2$, a finite uniform factor
relative to $g_M(k)=\max\{M^{-1/2},k/M\}$ exists if and only if
$\liminf_{M\to\infty}r_M/M>0$, for stationary or arbitrary starts.
For each fixed $0<\alpha\le1/4096$ and
$r_M=\max\{2,\lceil\alpha M\rceil\}$,
\[
 2^{-26}\log\frac1{32\alpha}
 \le C_{\rm stat}(\alpha)
 \le C_{\rm all}(\alpha)
 \le2^{19}\left(1+\log\frac2\alpha\right).
\]
With exact counts and protected terminal pair, optionally also the
initial pair, a post-corruption full-word encoder needs at least two
supplemental symbols to attain a finite uniform factor; two suffice
when its message is exact and fault-free.
\end{proposition}
\begin{proof}
Suppose first that the liminf is zero. Select an increasing subsequence
on which $r_M/M\to0$, and on that subsequence put
\[
 b_M=\max\{r_M,\sqrt M\},\qquad
 k_M=\lfloor\sqrt{Mb_M}\rfloor.
\]
Then $b_M=o(M)$, $k_M/M\to0$, $k_M/b_M\to\infty$,
$r_M/k_M\to0$, and $k_M/\sqrt M\to\infty$. Moreover
$M/k_M^2\to0$, so every condition in
\eqref{eq:finite-lower-domain} holds eventually. Here
$t_M=ck_M/M$, $\rho=u_0/(8Mt_M)$, $L_-\to\infty$, and
$\rho\to1/32$.
If one $k$-blind family had a finite admissible factor $C$, its clean
risk on every alternative at this same $M$ and supplied certificate
would satisfy $S_M\le C/\sqrt M=o(t_M)$. For example, eventually
$\rho\ge1/64$ and $t_M/16-S_M\ge t_M/32$, so
\[
 \frac{R_{\rm null}}{g_M(k_M)}
 \ge\frac{c}{2048}L_-\longrightarrow\infty,
\]
since $g_M(k_M)=k_M/M$ eventually. This contradicts $C$.
The alternatives may depend on the tested budget, but the predictor
and its clean guarantee are the same map at the same horizon. These
stationary witnesses also rule out an all-start guarantee.

Conversely, if the liminf is positive, choose $\delta>0$ and $M_0$
so $r_M\ge\delta M$ for $M\ge M_0$. Put
$m_M=\min\{r_M-2,\lfloor M/2\rfloor\}$.
Then $(m_M+1)/M$ is bounded away from zero for all sufficiently large
$M$. The suffix estimate of Proposition~\ref{prop:finite-upper} is at
most $2^{19}[1+\log(M/(m_M+1))]g_M(k)$ at $\beta=1/M$, so gives
one uniform all-start factor eventually. A finite initial set of
horizons is harmless: any prediction has risk at most one and
$g_M(k)\ge M^{-1/2}$. This proves the exact criterion without
monotonicity of the widths.

Now fix $\alpha>0$ in the stated interval before sending $M$ to
infinity. For any family with finite admissible stationary factor $C$,
put $k=128r_M$ at sufficiently large horizons. Then
\[
 \frac kM\to128\alpha\le\frac1{32},\quad
 t\to128c\alpha>0,\quad At\to8\alpha<\frac14,
\]
all finite-domain conditions eventually hold, and
\[
 L_-\to\log\frac1{32\alpha},\qquad
 \frac{u_0}{8Mt}\to\frac1{32},\qquad
 S/t\le\frac{C}{t\sqrt M}\to0.
\]
For these budgets $k\ge\sqrt M$ eventually. Using the deliberately
loose $g_M(k)\le2k/M$ in \eqref{eq:finite-lower}, the positive limiting
bracket gives
\[
 C\ge\liminf_{M\to\infty}\frac{R_{\rm null}}{g_M(k)}
 \ge\frac1{32}\frac{c}{32}\log\frac1{32\alpha}
 =2^{-26}\log\frac1{32\alpha}.
\]
This holds for every admissible finite factor, so taking the infimum
does not require an attained optimum. The upper follows from
Proposition~\ref{prop:finite-upper}: $m_M+1\ge\alpha M/2$, including
$r_M=2$, hence $1+\log(M/(m_M+1))\le1+\log(2/\alpha)$.
The stationary class is a subclass of the arbitrary-start class.
These arguments take the horizon limit for each fixed fraction and
factor; they neither exchange limits nor give an every-budget
logarithmic risk lower bound.

Finally, a one-symbol supplemental alphabet is constant, so leaves
only counts and the supplied endpoint pair or pairs. At $r_M=2$ even
both endpoint pairs are insufficient by the necessity just proved.
The exact one-bit factorization in
Proposition~\ref{prop:finite-upper} supplies two symbols and has
all-start risk at most $(2^{17}+2^{23})g_M(k)$. This concerns a
fault-free message computed by scanning the whole already corrupted
word, not one raw coordinate or a message subject to further corruption.
\end{proof}

\subsection{Elementary same-information oracle lower floors}

\begin{lemma}[Critical oracle lower floors]
\label{lem:oracle-floor}
For every $M\ge64$ and $0\le k\le M$, the stationary minimax risk
with supplied $\beta=1/M$ is at least $g_M(k)/96$, even if the predictor
knows $k$ and observes the entire corrupted word. Consequently this
lower bound holds for each restricted observation considered here and
for the arbitrary-start minimax risk.
\end{lemma}
\begin{proof}
We give separate clean and edit witnesses; taking their maximum is
permitted in a minimax lower bound.

\paragraph{Clean floor.}
Put $N=M+2$ and use stationary IID Bernoulli sources
\[
 p_\pm=\tfrac12\pm\frac1{16\sqrt N}.
\]
Their pair chains satisfy $P^2(u,v)=\Pi(v)$ for every row, hence the
supplied certificate. The target is $p_\pm$ at every terminal pair.
The elementary Bernoulli divergence bound
\[
 d(p\Vert q)\le\frac{(p-q)^2}{q(1-q)}
\]
follows by applying $\log x\le x-1$ to its two summands.
With $\Delta=p_+-p_-=1/(8\sqrt N)$, it gives one-letter divergence
at most $1/(8N)$. Independence makes the full-word divergence at most
$1/8$, and their total variation distance is at most $1/4$.
For clarity, this last inequality follows without a statistical
independence assertion about a selected row: for finite probability
laws $P,Q$, group their masses on
$A=\{x:P(x)\ge Q(x)\}$ and its complement. The log-sum inequality
(Jensen's inequality for $x\log x$) gives
$D(P\Vert Q)\ge d(P(A)\Vert Q(A))$.
For fixed $q$, the second derivative in $p$ of $d(p\Vert q)$ is
$1/[p(1-p)]\ge4$; its value and first derivative vanish at $p=q$.
Thus $D(P\Vert Q)\ge2(P(A)-Q(A))^2=2\TV(P,Q)^2$.

Testing by the nearer target makes each error incur prediction loss
at least $\Delta/2$. Every possibly randomized test has average error
at least $(1-\TV(P,Q))/2$, since the difference of its acceptance
probabilities is at most total variation. Thus the average prediction
risk is at least
\[
 \frac\Delta4(1-\TV(P,Q))
 \ge\frac3{128\sqrt N}\ge\frac1{64\sqrt M}.
\]
Identity corruption is allowed at every $k$, including for a
known-budget predictor.

\paragraph{Edit floor.}
For $k\ge1$, use the uniform prior on IID Bernoulli parameters
$p\in[1/4,3/4]$. Put $b=k+1$. If the first $M$ symbols contain $n$
ones, retain a uniformly chosen set of
$r=b\lfloor n/b\rfloor$ of those ones, turning the others to zero.
Leave the final two symbols $U$ unchanged. This single kernel is
source-independent and changes $n-r\le b-1=k$ editable bits on every
outcome.
For a specified retained set $T$ of size $r$, conditional on $n$,
\[
 \Pr(T\text{ retained}\mid n)
 =\frac{\binom{M-r}{n-r}}
        {\binom Mn\binom nr}
 =\frac1{\binom Mr}.
\]
Hence conditional on $r,U$ the output arrangement is uniform and
independent of $p$, even after mixing the possible $n$ in a bucket.
The full observation factors through at most
\[
 D=4\bigl(\lfloor M/b\rfloor+1\bigr)
\]
labels $(r,U)$ followed by a $p$-independent channel. The four protected
pair possibilities are counted, not conditioned away.

Realize that channel using a seed independent of $p$ and of the label,
and likewise use an independent decoder seed. Freeze both seeds.
There are at most $D$ possible prediction values. For the uniform
prior on an interval of length $\ell=1/2$, the union of radius-$v$
neighborhoods of those values has prior mass at most $2Dv/\ell$,
regardless of how a label is randomly selected conditional on $p$.
Outside that union every possible prediction has error greater than
$v$. Integrating this probability bound gives Bayes absolute-error
risk at least
\[
 \int_0^{\ell/(2D)}\left(1-\frac{2Dv}{\ell}\right)dv
 =\frac\ell{4D}.
\]
Averaging the independent seeds preserves it. Consequently some
parameter in the prior support has risk at least
\[
 \frac1{32(\lfloor M/(k+1)\rfloor+1)}
 \ge\frac{k+1}{32(M+k+1)}\ge\frac{k}{96M}.
\]
IID stationarity and $P^2(u,v)=\Pi(v)$ again give the same source
class and supplied certificate, with target $p$ at every row. Combining
the clean floor and edit floor yields $g_M(k)/96$, also at $k=0$.
A predictor with less information can only have a larger minimax risk,
and allowing arbitrary starts contains the stationary experiment.
\end{proof}
